% GENERATED FILE. Edit canonical lessons, then run npm run generate:books.
% canonical-source-hash: fnv1a64:c556aeb674a3a871
% canonical-lessons: ZH-C240-shengbing, ZH-C240-fayan, ZH-C240-zhong, ZH-C240-tu, ZH-C240-laduzi

\chapter{Feeling Ill}
\label{ch:zh-feeling-ill}
\hlchaptermodality{240}

\begin{chapteropening}
\textbf{By the end of this chapter:} \emph{I can say I'm ill and why: an infection, a swelling, being sick and an upset stomach.}

Complete the last lesson of chapter 240: I can say I'm ill and why: an infection, a swelling, being sick and an upset stomach.
\end{chapteropening}

\section[\texorpdfstring{\zh{生病} (shēngbìng)}{生病 (shēngbìng)}]{\zh{生病} (shēngbìng) — to get sick; to fall ill}

\label{lesson:ZH-C240-shengbing}



\begin{quote}
Before the new one: say the Mandarin for to breathe, then the Mandarin for a medical check-up.
\end{quote}

\subsection*{You'll want to know: \zh{生病}}
\textbf{\zh{生病}} — \emph{shēngbìng} — \textquotedblleft{}to get sick; to fall ill\textquotedblright{}.

\begin{quote}
\textbf{\zh{我生病了头很疼不想吃东西}} \emph{wǒ shēngbìng le tóu hěn téng bù xiǎng chī dōngxi} — I'm ill: my head really hurts and I don't want to eat.
\end{quote}

\begin{cousinweb}
\textbf{\zh{生病}} (\emph{shēngbìng}) is built from \zh{生} 'to grow' and \zh{病} 'illness', both inherited from Middle Chinese, \emph{srjaeng} and \emph{bjaengH} (Guangyun, in Baxter's transcription), which glosses the second 'worry' and 'suffering' and explains it, after an older book, as an illness grown worse.
\end{cousinweb}

\begin{grammarlens}[title={What you've built}]
One more word for giving reasons.
\end{grammarlens}

\subsection*{Your turn}
\begin{itemize}
\raggedright
  \item \emph{Say it:} \emph{shēngbìng}
  \item \emph{Say it:} \emph{shēngbìng}, once more
  \item \emph{Say it:} \emph{wǒ shēngbìng le tóu hěn téng bù xiǎng chī dōngxi} — I'm ill: my head really hurts and I don't want to eat.
  \item \emph{Recall:} say the Mandarin for to breathe, then the Mandarin for a medical check-up, then say \emph{shēngbìng} again
\end{itemize}

\begin{tcolorbox}[breakable,colback=teal!4,colframe=teal!35!black,title={Before you move on}]
What does \zh{生病} mean? (\textquotedblleft{}To get sick; to fall ill\textquotedblright{}.) Say it once more.
\end{tcolorbox}
\section[fāyán]{fāyán — to become inflamed}

\label{lesson:ZH-C240-fayan}



\begin{quote}
Before the new one: say the Mandarin for a medical check-up, then the Mandarin for to get sick.
\end{quote}

\subsection*{You'll want to know: fāyán}
\emph{fāyán} — \textquotedblleft{}to become inflamed\textquotedblright{}.

\begin{quote}
\emph{wǒ de sǎngzi fāyán le} — My throat is inflamed.
\end{quote}

Its characters come later; for now it is a word for the ear.

\begin{cousinweb}
\emph{fāyán} is built from \emph{fā} 'to start' and \emph{yán} 'flame', both inherited from Middle Chinese, \emph{pjot} and \emph{hjem} (Guangyun, in Baxter's transcription), which glosses the second 'heat' and explains it, after an older book, as fire rising.
\end{cousinweb}

\begin{grammarlens}[title={What you've built}]
One more word for giving reasons.
\end{grammarlens}

\subsection*{Your turn}
\begin{itemize}
\raggedright
  \item \emph{Say it:} \emph{fāyán}
  \item \emph{Say it:} \emph{fāyán}, once more
  \item \emph{Say it:} \emph{wǒ de sǎngzi fāyán le} — My throat is inflamed.
  \item \emph{Recall:} say the Mandarin for a medical check-up, then the Mandarin for to get sick, then say \emph{fāyán} again
\end{itemize}

\begin{tcolorbox}[breakable,colback=teal!4,colframe=teal!35!black,title={Before you move on}]
What does fāyán mean? (\textquotedblleft{}To become inflamed\textquotedblright{}.) Say it once more.
\end{tcolorbox}
\section[zhǒng]{zhǒng — swollen}

\label{lesson:ZH-C240-zhong}



\begin{quote}
Before the new one: say the Mandarin for to get sick, then the Mandarin for to become inflamed.
\end{quote}

\subsection*{You'll want to know: zhǒng}
\emph{zhǒng} — \textquotedblleft{}swollen\textquotedblright{}.

\begin{quote}
\emph{tā de jiǎo zhǒng le} — His foot has swollen up.
\end{quote}

Its characters come later; for now it is a word for the ear.

\begin{cousinweb}
\emph{zhǒng} is inherited from Middle Chinese \emph{tsyowngX} (Guangyun, in Baxter's transcription), glossed 'an illness' and, after an older book, 'a boil'.
\end{cousinweb}

\begin{grammarlens}[title={What you've built}]
One more word for giving reasons.
\end{grammarlens}

\subsection*{Your turn}
\begin{itemize}
\raggedright
  \item \emph{Say it:} \emph{zhǒng}
  \item \emph{Say it:} \emph{zhǒng}, once more
  \item \emph{Say it:} \emph{tā de jiǎo zhǒng le} — His foot has swollen up.
  \item \emph{Recall:} say the Mandarin for to get sick, then the Mandarin for to become inflamed, then say \emph{zhǒng} again
\end{itemize}

\begin{tcolorbox}[breakable,colback=teal!4,colframe=teal!35!black,title={Before you move on}]
What does zhǒng mean? (\textquotedblleft{}Swollen\textquotedblright{}.) Say it once more.
\end{tcolorbox}
\section[tù]{tù — to throw up; to vomit}

\label{lesson:ZH-C240-tu}



\begin{quote}
Before the new one: say the Mandarin for to become inflamed, then the Mandarin for swollen.
\end{quote}

\subsection*{You'll want to know: tù}
\emph{tù} — \textquotedblleft{}to throw up; to vomit\textquotedblright{}.

\begin{quote}
\emph{wǒ yùnchuán yǒudiǎnr xiǎng tù} — I'm seasick and feel a bit like throwing up.
\end{quote}

Its characters come later; for now it is a word for the ear.

\begin{cousinweb}
\emph{tù} is inherited from Middle Chinese \emph{thuH} (Guangyun, in Baxter's transcription), glossed 'to vomit'; the dictionary also lists the character with a rising tone for spitting out.
\end{cousinweb}

\begin{grammarlens}[title={What you've built}]
One more word for giving reasons.
\end{grammarlens}

\subsection*{Your turn}
\begin{itemize}
\raggedright
  \item \emph{Say it:} \emph{tù}
  \item \emph{Say it:} \emph{tù}, once more
  \item \emph{Say it:} \emph{wǒ yùnchuán yǒudiǎnr xiǎng tù} — I'm seasick and feel a bit like throwing up.
  \item \emph{Recall:} say the Mandarin for to become inflamed, then the Mandarin for swollen, then say \emph{tù} again
\end{itemize}

\begin{tcolorbox}[breakable,colback=teal!4,colframe=teal!35!black,title={Before you move on}]
What does tù mean? (\textquotedblleft{}To throw up; to vomit\textquotedblright{}.) Say it once more.
\end{tcolorbox}
\section[lā dùzi]{lā dùzi — to have diarrhoea; to have an upset stomach}

\label{lesson:ZH-C240-laduzi}



\begin{quote}
Before the new one: say the Mandarin for swollen, then the Mandarin for to throw up.
\end{quote}

\subsection*{You'll want to know: lā dùzi}
\emph{lā dùzi} — \textquotedblleft{}to have diarrhoea; to have an upset stomach\textquotedblright{}.

\begin{quote}
\emph{tā zuótiān lā dùzi le} — He had an upset stomach yesterday.
\end{quote}

Its characters come later; for now it is a word for the ear.

\begin{cousinweb}
\emph{lā dùzi} is built from \emph{lā} 'to pull' and \emph{dùzi} 'belly', all three syllables inherited from Middle Chinese, \emph{lop}, \emph{duX} and \emph{tsiX} (Guangyun, in Baxter's transcription), which glosses the first 'to break'; how the phrase came to mean an upset stomach is not settled.
\end{cousinweb}

\begin{grammarlens}[title={What you've built}]
One more word for giving reasons.
\end{grammarlens}

\subsection*{Your turn}
\begin{itemize}
\raggedright
  \item \emph{Say it:} \emph{lā dùzi}
  \item \emph{Say it:} \emph{lā dùzi}, once more
  \item \emph{Say it:} \emph{tā zuótiān lā dùzi le} — He had an upset stomach yesterday.
  \item \emph{Recall:} say the Mandarin for swollen, then the Mandarin for to throw up, then say \emph{lā dùzi} again
\end{itemize}

\begin{tcolorbox}[breakable,colback=teal!4,colframe=teal!35!black,title={Before you move on}]
What does lā dùzi mean? (\textquotedblleft{}To have diarrhoea; to have an upset stomach\textquotedblright{}.) Say it once more.
\end{tcolorbox}
